\section{Experiment Skeleton}

\subsection{Purpose}

The experiments in this project are verification experiments rather than
systems-performance experiments. Their role is to validate the specification,
surface counterexamples early, and document the cost and scope of the bounded
checks used during development.

\subsection{Experiment 1: Small-Instance Model Checking}

This experiment explores tiny bounded instances of the state space with a small
account set, bounded prices, and bounded quantities. The goal is to catch
specification bugs before attempting full invariant proofs.

\begin{itemize}[leftmargin=1.5em]
\item \textbf{Inputs:} number of accounts, maximum outstanding orders, bounded
price range, bounded quantity range.
\item \textbf{Checks:} stable-book violations, queue-discipline violations,
impossible trades, and quantity inconsistencies.
\item \textbf{Outputs to report:} maximum bounds explored, whether a
counterexample was found, and the kind of bug revealed by any failing case.
\end{itemize}

\subsection{Experiment 2: Invariant Checking}

This experiment focuses on proving the main safety properties with Veil's
invariant tooling.

\begin{itemize}[leftmargin=1.5em]
\item \textbf{Goal:} determine whether the chosen invariants are strong enough
to make the main properties inductive.
\item \textbf{Measurements:} which invariants are required, whether
counterexamples to induction appear, and which helper lemmas are needed from
Lean.
\item \textbf{Success criterion:} the core safety set passes invariant
checking for the minimal model.
\end{itemize}

\subsection{Experiment 3: Optional Executable Cross-Check}

If a Velvet matching loop is implemented, compare its outputs on bounded event
streams against the abstract Veil postcondition.

\begin{itemize}[leftmargin=1.5em]
\item \textbf{Goal:} show that the executable loop agrees with the abstract
matching relation on representative small traces.
\item \textbf{Outputs to report:} supported trace sizes, observed mismatches if
any, and the final postcondition used to link code back to the model.
\end{itemize}

\subsection{Reporting Template}

\begin{center}
\begin{tabular}{@{}>{\raggedright}p{3.6cm} p{4.2cm} p{3.0cm} p{4.0cm}@{}}
\textbf{Experiment} & \textbf{Purpose} & \textbf{Variables} & \textbf{What to report} \\
\hline
Bounded model checking & Find specification bugs early & Accounts, prices, quantities, outstanding orders & Largest bounds tried; counterexamples found or absent \\
Invariant checking & Prove core safety properties & Candidate invariants, helper lemmas & Which properties proved; missing obligations \\
Optional executable check & Align code with spec & Trace length, supported transition subset & Agreement or mismatch between loop and spec \\
\end{tabular}
\end{center}
